\subsection{Electronic reference and state assignment}
\label{sec:electronic-reference}

The automatic gas-phase workflow uses an all-electron, closed-shell
restricted Kohn--Sham (RKS) reference and real, spin-adapted singlet states
in the Tamm--Dancoff approximation (TDA) to time-dependent density
functional theory.\cite{pyscf,tda} The default is PBE0/def2-TZVPD with an
unpruned atom-centered integration grid.\cite{pbe0,diffuse} At Cartesian
geometry $R$ the energy of the selected excited surface is
\begin{equation}
 A(R) q = \omega q, \qquad q^{\mathsf T}q=1, \qquad
 E_{\mathrm{ex}}(R)=E_0(R)+\omega(R).
 \label{eq:electronic-tda}
\end{equation}
Here $E_0$ is the total restricted ground-state energy and $A$ is the real
symmetric TDA operator in the occupied--virtual orbital space. PySCF
stores the excitation amplitudes as $X=q/\sqrt{2}$, with zero
de-excitation amplitudes. The transpose in Eq.~\eqref{eq:electronic-tda}
denotes the Euclidean scalar product after flattening the
occupied--virtual indices.

At each geometry the CPU workflow requires native self-consistent field
(SCF) convergence and a norm of the occupied--virtual orbital gradient
of at most $10^{-7}\,E_h$ by default, evaluated with a fresh Fock matrix
from the final density. It tries a projected occupied space from a nearby
geometry when available, followed by fresh restricted SCF guesses with
first- and second-order solvers. The requested TDA eigenpairs must be
converged, positive and sorted, with $2X_k^{\mathsf T}X_l=\delta_{kl}$ in
the flattened representation and $\lVert A q_k-\omega_k q_k\rVert_2\le
10^{-8}\,E_h$ by default. If the Davidson solution fails validation, a
small particle--hole space can be diagonalized explicitly using the same
operator. Larger spaces use a restarted iterative solve. The excited-state
gradient also checks the absolute norm of the orbital-gap-preconditioned
coupled-perturbed Hartree--Fock (CPHF) Z-vector residual against the
response tolerance ($10^{-10}$ by default). A residual correction uses
the original response operator and tolerance.

The user selects a root at the ground-state minimum. At the next
geometry $R'$, the occupied and virtual blocks $O$ and $V$ of
\begin{equation}
 C(R)^{\mathsf T} S(R,R') C(R'), \qquad
 S_{\mu\nu}(R,R')=\langle\phi_\mu(R)\mid\phi_\nu(R')\rangle
\end{equation}
transport the two indices of the amplitude. Here $C$ contains the
molecular orbital coefficients and $\phi_\mu$ is an atomic orbital.
The signed score of root $k$ at $R'$ is
\begin{equation}
 s_k =
 \frac{\left\langle O^{\mathsf T}X_{\rm ref}V, X_k(R')\right\rangle_F}
 {\lVert X_{\rm ref}\rVert_F\,\lVert X_k(R')\rVert_F},
 \label{eq:electronic-tracking}
\end{equation}
where $\langle\cdot,\cdot\rangle_F$ is the Frobenius scalar product and
$X_{\rm ref}$ carries the phase of the accepted preceding point. This
overlap of particle--hole transition densities is normalized by the two
original norms so that it retains any loss on transport. The workflow
maximizes $|s_k|$ and multiplies the selected $X_k$ and its dipole by
$\operatorname{sgn}(s_k)$. The minimum score magnitude is 0.65. With
multiple roots, the automatic workflow also requires a margin of at
least 0.10 over the runner-up and a gap of at least $10^{-4}\,E_h$ to
the closest computed neighbor. It
expands the root window on insufficient overlap or a match at its upper
edge, within a bounded search. A trial that changes the energy rank of
the tracked root is rejected in the isolated adiabatic model.

\subsection{Stationary geometries and local derivative fits}
\label{sec:electronic-geometry}

The gas-phase adiabatic Hessian (AH) workflow optimizes both surfaces
with analytic gradients. It projects translations and rotations out
of the Cartesian Broyden--Fletcher--Goldfarb--Shanno (BFGS) step and
initializes the excited-state inverse Hessian from the projected
ground-state Hessian. Each search starts with a maximum per-atom step
length of at most 0.22~bohr before alignment and backtracking. With $g$
the projected gradient and $\Delta R$ the realized displacement after
alignment, the step must pass electronic validation and satisfy
$g^{\mathsf T}\Delta R<0$ and
\[
 E(R+\Delta R)\le E(R)+10^{-4}g^{\mathsf T}\Delta R+2\times10^{-11}\,E_h .
\]
Steps that fail are shortened by successive halving. If the
quasi-Newton search fails, the optimizer can restart from the last
accepted geometry in the steepest-descent direction. The maximum norm
of the Cartesian gradient on any atom must be at most
$10^{-5}\,E_h\,\mathrm{bohr}^{-1}$ at a stationary endpoint. The
minimum check requires positive, resolved curvature in the vibrational
subspace. For a clearly negative eigenvalue the workflow searches for
downhill branches along both signs of the lowest-curvature mode.

The restricted determinant at each stationary geometry is tested for
orbital instabilities within the restricted space and toward an
unrestricted solution.\cite{stability} These tests probe orbital
curvature in addition to SCF stationarity. An internal instability can
restart the workflow with orbital-stability checks at every geometry.
An unstable orbital rotation then seeds a fresh restricted SCF solve
that must lower the energy and satisfy the same residual criterion.
A verified instability toward an unrestricted determinant stops the
automatic calculation with that functional. Only explicitly supplied
alternate functionals are eligible for a fresh attempt, with the
target root assigned anew at the corresponding ground-state minimum.
An inconclusive stability solve terminates preparation.

The automatic RKS/TDA route uses the analytic PySCF ground-state
Hessian at fixed quadrature and the native TD-RKS excited-state
gradient without an explicit moving-grid correction. It estimates
excited-state curvature by fitting $g(x)-g(0)\simeq H x$ in a rigid-motion
free Cartesian subspace, with $H=H^{\mathsf T}$ and $x$ a displacement
in that subspace. Independent displacements test both the predicted
gradients and the quadratic energy expansion. First-order
Herzberg--Teller (HT) slopes are fitted to signed dipole differences
in ground-state normal coordinates and checked on independent points.
Each fit can be attempted at successively smaller radii. For HT, when
the vibrational quantum of the dominant dipole-active mode reaches
the neighboring electronic gap and its predicted dipole change is
resolvable, direct calculations at both signed zero-point displacements
also test the linear expansion. The workflow requires these checks to
pass before using the derivative model. The explicitly selected
vertical-gradient (VG) model instead uses the ground-state Hessian and
the excited-state gradient at the ground-state minimum to construct
two wells of equal curvature with a constant transition dipole.

\subsection{Analytic derivatives of prepared TDA states}
\label{sec:electronic-analytic}

The prepared-state GPU4PySCF route differentiates isolated real singlet
TDA roots of all-electron, global-hybrid LDA or GGA references in
equilibrium linear-response conductor-like polarizable continuum
solvation (LR C-PCM).\cite{gpu,cpcm,liangtda,liangpcm} Its curvature
implementation uses spherical atomic orbitals, Coulomb-metric density
fitting for both Coulomb and exchange, and the same switching/Gaussian
(SWIG) cavity for the SCF and TDA calculations. It forms derivatives
in an orthonormal moving orbital frame. In that frame the action on
an occupied--virtual matrix $U$ is
\begin{equation}
 A[U]=UF_{vv}-F_{oo}U
       +2C_o^{\mathsf T}\mathcal{V}[C_v U^{\mathsf T}C_o^{\mathsf T}]
                 ^{\mathsf T}C_v ,
\end{equation}
where $F_{oo}$ and $F_{vv}$ are the occupied and virtual Fock blocks
and $C_o$ and $C_v$ are the corresponding orbital coefficients.
The linear singlet response $\mathcal{V}$ acts on the unsymmetrized
atomic-orbital transition density and includes Coulomb, exchange,
semilocal exchange--correlation, and reaction-field terms.
The derivative implementation checks this action on the target root
against the native operator. With Cartesian derivative
indices $\alpha,\beta$, define
$P=I-qq^{\mathsf T}$, $L=P(\omega I-A)P$, and
$A_\alpha=\partial_\alpha A$, including the response of the RKS
reference. On the subspace orthogonal to $q$, differentiation of
Eq.~\eqref{eq:electronic-tda} gives
\begin{align}
 \omega_\alpha &= q^{\mathsf T} A_\alpha q, &
 L q_\alpha &= P A_\alpha q, &
 q^{\mathsf T} q_\alpha &=0, \\
 H^{\mathrm{ex}}_{\alpha\beta}
 &= H^{0}_{\alpha\beta}
 + q^{\mathsf T} A_{\alpha\beta}q
 + 2(PA_\alpha q)^{\mathsf T} q_\beta .
 \label{eq:electronic-hessian}
\end{align}
The native ground-state Hessian $H^0$ includes motion of the
quadrature and auxiliary basis. The mixed contraction at fixed $q$
contains first-order occupied orbital responses and an orbital
adjoint that eliminates the unknown second-order occupied--virtual
rotation. The integral terms include motion of the orbital and
auxiliary bases in the three-center integrals and inverse Coulomb
metric, with $\partial_\alpha M^{-1}=-M^{-1}M_\alpha M^{-1}$ for
the auxiliary metric $M$. The solvent terms include derivatives of
the ground-state reaction field and of the response to the
transition density. The semilocal exchange--correlation terms
differentiate both the grid nodes and weights and include functional
derivatives through fourth order.

The projected solves check the eigenpair residual and separation of
the target from neighboring roots and verify each response residual.
Curvature uses a matrix-free minimum-residual solve, with batched
conjugate-gradient solves for the lowest root when available and
individual minimum-residual solves for unresolved rows. The three
contributions to Eq.~\eqref{eq:electronic-hessian} are checked separately
for symmetry before averaging numerical asymmetry. The TDA response
alone is solved once more at tighter tolerances if its contraction
exceeds $10^{-8}\,E_h\,\mathrm{bohr}^{-2}$ in asymmetry. A failed
check ends the prepared analytic calculation.

For a Cartesian dipole component $d$, let $b_d$ be the vector of
$\sqrt{2}\langle i|r_d|a\rangle$ in the native length-gauge sign
convention. Solving an adjoint equation for each Cartesian component gives
\begin{equation}
 \mu_d=b_d^{\mathsf T}q, \qquad
 L z_d=P b_d, \qquad
 \partial_\alpha\mu_d
 =(\partial_\alpha b_d)^{\mathsf T}q+z_d^{\mathsf T} A_\alpha q.
 \label{eq:electronic-dipole}
\end{equation}
The first contraction includes analytic position-integral derivatives
and occupied and virtual orbital motion. The second uses polarization
of quadratic TDA gradient contractions at the same geometry, with a
moving exchange--correlation quadrature correction. A selected sign
multiplies both $\boldsymbol\mu$ and its derivative. A separate
conventional RHF configuration-interaction singles (CIS) implementation
supplies dense analytic Hessians and dipole derivatives for pure HF.
The automatic pure-HF route can fall back to the local fits when these
dense calculations exceed the memory limit. Analytic derivatives
that fail validation terminate the calculation.

\subsection{Spin-flip optical response}

The optional spin-flip route starts from a converged high-spin
unrestricted Kohn--Sham triplet. Its multi-collinear TDA calculation
produces spin-flip-down roots with energies $\Omega_i$ relative to that
auxiliary reference.\cite{gpu} The optical gap between two selected
low-spin roots is $\Delta_{ji}=\Omega_j-\Omega_i$. For their real amplitudes
$X_i$ and $X_j$, the occupied and virtual blocks of the mixed density
have the atomic-orbital form
\begin{equation}
 D_{\alpha}^{\mathrm d}=-C_o^\alpha X_i X_j^{\mathsf T}(C_o^\alpha)^{\mathsf T},
 \qquad
 D_{\beta}^{\mathrm d}=C_v^\beta X_j^{\mathsf T}X_i(C_v^\beta)^{\mathsf T}.
 \label{eq:sf-density}
\end{equation}
Orbital relaxation supplies the occupied--virtual parts. With
$\mathcal L_{\mathrm{sc}}$ the spin-conserving TDDFT action on the
excitation and de-excitation blocks, the implementation solves
$(\mathcal L_{\mathrm{sc}}-\Delta_{ji} I)z=b_{ij}$ by GMRES and verifies
its residual. The source $b_{ij}$ contains the response to the blocks in
Eq.~\eqref{eq:sf-density}, third-order exchange--correlation derivatives,
and commutators of the transverse potentials with the spin-flip transition
densities. Contracting the relaxed mixed density with the Cartesian
position integrals gives the length-gauge dipole. This is an optical
transition between the selected roots. The population and photon-energy
weights enter only in the subsequent vibronic assembly.

\subsection{Full-space spin--orbit response and analytic Herzberg--Teller coupling}
\label{sec:electronic-phosphorescence}

The phosphorescence model uses first-order spin--orbit coupling (SOC)
from the molecular Breit--Pauli spin--orbit mean-field (SOMF) operator
on an all-electron RKS reference at the emitting triplet minimum.
\cite{somf,desouza-phosphorescence} Let $A_S$ and $A_T$ be the
spin-adapted singlet and triplet TDA operators in the same spatial
occupied--virtual basis. The target triplet has unit amplitude $Q$
and excitation energy $E_T$ relative to the RKS ground state. In a
Cartesian triplet spin basis, $h_m$ is the imaginary part of the
spatial SOC coefficient multiplying $\sigma_m/2$, where $\sigma_m$ is
a Pauli matrix.
For a one-electron operator $O$, define its particle and hole action
by $\ell_O Q=QO_{vv}^{\mathsf T}-O_{oo}^{\mathsf T}Q$. For
Cartesian component $\xi$ of the physical electric dipole $\mu_\xi$,
use $d_\xi=\sqrt{2}(\mu_\xi)_{ov}$, $g_m=(h_m)_{ov}/\sqrt{2}$,
$b_m=\ell_{h_m}Q/2$, and $t_\xi=\ell_{\mu_\xi}Q$.
With all occupied--virtual arrays flattened in the contractions,
\begin{align}
 B_S x_\xi &= d_\xi, & B_S &= A_S-E_T I, &
 A_T y_\xi &= t_\xi, \\
 M_{m\xi} &= -i\bigl(b_m^{\mathsf T}x_\xi+g_m^{\mathsf T}y_\xi\bigr).
 \label{eq:full-space-soc}
\end{align}
The two amplitudes describe singlet admixture into the emitting triplet
and triplet admixture into the ground singlet. The common permanent
dipole cancels between the two state responses. The shifted operators
must be positive, and the response solves span the full finite
occupied--virtual space without explicit enumeration of its roots.

For a mass-weighted triplet normal coordinate $s$, differentiating the
response equations gives the first-order HT coefficient
\begin{align}
 \partial_s M_{m\xi}=-i\bigl[&
 (\partial_s b_m)^{\mathsf T}x_\xi
 +u_m^{\mathsf T}\bigl(\partial_s d_\xi-(\partial_s B_S)x_\xi\bigr)
 \notag\\
 &+(\partial_s g_m)^{\mathsf T}y_\xi
 +v_m^{\mathsf T}\bigl(\partial_s t_\xi-(\partial_s A_T)y_\xi\bigr)
 \bigr],
 \label{eq:soc-ht}
\end{align}
where $B_S u_m=b_m$ and $A_T v_m=g_m$. The derivatives include the
response of the SOMF and position integrals, the spin-summed density,
the orbitals, the two TDA operators, and the isolated spatial triplet
root.\cite{pygraso} We evaluate these terms analytically at one gas-phase
geometry, without fitting or finite differences. The first-order SOC
model neglects the fine-structure splittings and uses equal spin
populations on a common triplet vibrational surface. The nuclear
calculation adds the Condon and linear HT amplitudes coherently for each
spin component before forming the intensity. The rate uses the vacuum
photon prefactor and weights each positive photon energy by its cube.
